\chapter{Related Work}

\section{Traffic Sign Classification using K-d trees and Random Forests (Zaklouta et al., 2011)}
In this research, Zaklouta et al. (2011) evaluated the performance of K-d trees and Random Forests for traffic sign classification using the full 43-class GTSRB dataset. Their methodology focused on two primary feature extraction techniques: Histogram of Oriented Gradients (HOG) and Distance Transforms. The researchers found that HOG descriptors, particularly when computed on 40x40 pixel images, provided highly discriminative features for tree-based classifiers. Using an ensemble of 500 trees in a Random Forest, they achieved a peak classification accuracy of \textbf{97.2\%} \cite{zaklouta2011}. This study highlighted that Random Forests are less sensitive to background variations and imbalanced data compared to K-d trees, which achieved 92.9\% accuracy under similar conditions.

\section{Traffic sign detection and recognition based on random forests (Ellahyani et al., 2016)}
Ellahyani et al. (2016) proposed a recognition module that integrates both color and texture information to improve classification robustness. Their approach extended the standard HOG features into the HSI (Hue, Saturation, Intensity) color space, resulting in a 5292-dimensional HSI-HOG descriptor. This was further combined with Local Self-Similarity (LSS) features to capture complex structural patterns. The classification was performed using a Random Forest with 750 trees, which achieved a Correct Classification Rate (CCR) of \textbf{97.43\%} on the GTSRB dataset \cite{ellahyani2016}. Their benchmarking tests demonstrated that the Random Forest classifier outperformed SVM with an RBF kernel (which reached 96.91\% accuracy) while being significantly more efficient in terms of computational speed.

\section{Comparison and Research Context}
While both Zaklouta et al. (2011) and Ellahyani et al. (2016) achieved state-of-the-art results on the full 43-class GTSRB dataset utilizing traditional feature engineering (HOG, LSS) and Random Forests, our current study explores a representative 10-class subset. We leverage modern deep feature extractors (ResNet50 and InceptionV3) to automate the representation learning process. By achieving \textbf{97.24\%} accuracy on this subset using an InceptionV3-SVM hybrid, our work demonstrates that the discriminative power of deep manifolds is comparable to the rigorous handcrafted feature engineering presented in these highly cited studies.
